\documentclass[11pt]{article}
\usepackage[margin=1in]{geometry}
\usepackage{amsmath,amssymb,amsthm}
\usepackage{graphicx}
\usepackage{booktabs}
\usepackage{hyperref}
\usepackage{xcolor}
\usepackage{algorithm}
\usepackage{algorithmic}
\usepackage{float}
\usepackage{multirow}
\usepackage{caption}
\usepackage{subcaption}
\usepackage[numbers]{natbib}
\usepackage{url}

\hypersetup{colorlinks=true, linkcolor=blue!60!black, citecolor=green!50!black, urlcolor=blue!70!black}

\newcommand{\EN}{E_N}
\newcommand{\ETN}{E_{2N}}
\newcommand{\MTN}{M_{2N}}
\newcommand{\vfinal}{v_{\text{final}}}
\newcommand{\MRE}{\text{MRE}}

\title{\textbf{Two-Phase PageRank with Travel-Time Self-Loops\\ for Urban Traffic Distribution Estimation\\ and Event-Driven Anomaly Detection}}

\author{
  Research Team\\
  \textit{University of Utah}\\
  Salt Lake City, Utah, USA
}

\date{}

\begin{document}
\maketitle

% ================================================================
\begin{abstract}
We present a modified PageRank algorithm for estimating link-level traffic distributions on urban road networks using sparse GPS trajectory data.
The model introduces two key innovations: (1)~\emph{travel-time self-loops} that weight each link's mass retention proportionally to its traversal time $\mu \cdot \ell_i / s_i$, and (2)~a \emph{two-phase block structure} separating active driving from local dispersal.
Applied to the Salt Lake City road network (99,716 directed links, 467,911 GPS trajectories over 30~days), the model achieves a Top-100 Mean Relative Error (MRE) of 0.0312 and Pearson correlation $r = 0.997$ against smoothed ground truth.

We further introduce \emph{demand-weighted transitions}, a single parameter $\gamma$ that encodes average traffic demand into the transition matrix, enabling the damping factor to be raised from 0.80 to 0.96 (reducing ground-truth dependence from 20\% to 4\%) while \emph{improving} overall accuracy.
This is validated through temporal, day-of-week, and 5-fold cross-validation splits.

For anomaly detection, we demonstrate that the dwell-time parameter $\mu$ serves as an effective per-link control knob: reducing $\mu$ on disrupted links from 20 to 0.01 drops the crash-corridor prediction error from 179\% to 11.5\%, and game-day stadium-area error from 293\% to 7.8\%---all without using any event-day ground truth.
Cross-validated Tuesday baseline analysis confirms $\mu = 20$ is optimal on normal days ($r = 0.9999$ between data gap and model error), establishing that deviations from this default reliably signal real-world disruptions.
\end{abstract}

% ================================================================
\section{Introduction}
\label{sec:intro}

Accurate estimation of link-level traffic distributions is fundamental to transportation planning, real-time route guidance, and incident detection~\cite{wardrop1952road, jiang2009citywide}.
Traditional approaches rely on fixed sensor networks (loop detectors, cameras), which provide high temporal resolution but limited spatial coverage.
The proliferation of GPS-equipped mobile devices has created an alternative data source: sparse trajectory observations that cover the entire network but observe each link infrequently~\cite{herrera2010evaluation, zheng2014urban}.

The core challenge is converting sparse, noisy GPS traversal counts into a complete traffic distribution over all links.
In a typical 15-minute observation window, only 1--3\% of links in a metropolitan network have any GPS observations~\cite{horvitz2012prediction}.
Directly using these counts produces a distribution dominated by zeros, with no information about unobserved links.

We address this through a three-stage pipeline:
\begin{enumerate}
    \item \textbf{Graph-diffusion smoothing} propagates observed counts to neighbouring links using the road adjacency structure, filling gaps while preserving the spatial correlation of traffic patterns.
    \item \textbf{Two-Phase PageRank with travel-time self-loops} computes a stationary distribution that balances observed data (teleportation) against network topology (transition matrix), producing physically realistic traffic estimates.
    \item \textbf{Demand-weighted transitions} encode average traffic demand into the matrix, enabling high damping factors ($d > 0.90$) that reduce dependence on per-timeframe ground truth.
\end{enumerate}

The remainder of this paper is organised as follows.
Section~\ref{sec:data} describes the data pipeline.
Section~\ref{sec:model} presents the two-phase model.
Section~\ref{sec:demand} introduces demand-weighted transitions.
Section~\ref{sec:results} reports accuracy results.
Section~\ref{sec:events} demonstrates event detection and prediction.
Section~\ref{sec:mu} analyses the role of~$\mu$.
Section~\ref{sec:conclusion} concludes.


% ================================================================
\section{Data Pipeline}
\label{sec:data}

\subsection{Study Area and Network}

The study area is the greater Salt Lake City metropolitan region, represented as a directed graph $G = (V, E)$ with $|E| = N = 99{,}716$ directed links and average out-degree 2.87.
The network is derived from OpenStreetMap via SUMO, with each link carrying attributes: length~$\ell_i$ (metres), free-flow speed~$s_i$ (m/s), and lane count.
Link lengths range from 0.59\,m to 11,577\,m (mean 142\,m), and speeds from 2.2 to 33.5\,m/s.

\subsection{GPS Trajectory Data}

We use 467,911 GPS trajectories recorded between August~31 and September~30, 2018.
Each trajectory consists of timestamped $(latitude, longitude)$ points from mobile devices.
The raw data is processed through:
\begin{enumerate}
    \item \textbf{Cleaning}: Speed filtering ($v > 2$\,m/s), 5-minute gap splitting into separate trips.
    \item \textbf{Map matching}: Hidden Markov Model-based matching~\cite{newson2009hidden} using the Leuven map-matching library~\cite{meert2018hmm}, producing sequences of traversed link IDs with entry timestamps.
    \item \textbf{Temporal binning}: Counting traversals per link per 15-minute window, yielding a $2{,}880 \times 63{,}118$ sparse count matrix (2.06\% density).
\end{enumerate}

\subsection{The Sparsity Problem}
\label{sec:sparsity}

In any 15-minute window, only 1--3\% of network links have GPS observations.
The raw count vector $c_t$ for timeframe~$t$ has $c_t[i] = 0$ for $\sim$97\% of links.
Directly using $c_t / \|c_t\|_1$ as a traffic distribution assigns zero probability to most links, making the Mean Relative Error (MRE) undefined for unobserved links.

\subsection{Graph-Diffusion Smoothing}
\label{sec:smoothing}

\subsubsection{Why Smoothing Is Necessary}

The raw count vector $c_t$ for any 15-minute window has $c_t[i] = 0$ for approximately 97\% of links (Figure~\ref{fig:sparsity}).
A probability distribution with 97\% zeros is not meaningful for traffic analysis---it implies that 97\% of roads carry exactly zero traffic, which is clearly false.
Any analysis that requires values on all links (routing, equilibrium computation, anomaly detection) must impute the missing 97\%.
The question is not \emph{whether} to smooth, but \emph{how} and \emph{how much}.

\begin{figure}[H]
\centering
\includegraphics[width=\linewidth]{fig11_sparsity_histogram.png}
\caption{\textbf{GPS data sparsity.} Distribution of link coverage (fraction of 99,716 links with $\geq 1$ GPS observation) per 15-minute window across all 2,880 timeframes. The median coverage is approximately 2\%, meaning 97--98\% of the network is unobserved in any given window.  Without imputation, any network-wide analysis is impossible.}
\label{fig:sparsity}
\end{figure}

\subsubsection{The Smoothing Operator}

We use a one-step graph-diffusion filter~\cite{shuman2013emerging}:
\begin{equation}
    \boxed{c_{\text{smooth}} = \bigl[(1-\gamma)\,I + \gamma\,P_{\text{adj}}\bigr]\,\tilde{c}}
    \label{eq:smoothing}
\end{equation}
where $\tilde{c} = c + \alpha \cdot \mathbf{1}$ is the Laplace-smoothed count vector ($\alpha = 0.1$), $P_{\text{adj}}$ is the row-stochastic adjacency matrix, and $\gamma \in [0, 1]$ is the diffusion strength.
The choice of graph diffusion is principled: it uses \emph{only the road adjacency structure} to propagate information, encoding the physical prior that adjacent links carry correlated traffic.

\subsubsection{Spectral Interpretation}

$P_{\text{adj}}$ has eigenvalues $\lambda \in [-1, 1]$.
The filter $(1-\gamma)I + \gamma P_{\text{adj}}$ has eigenvalues $h(\lambda) = (1-\gamma) + \gamma\lambda$.
This acts as a \emph{low-pass filter} on the graph spectrum:
\begin{itemize}
    \item \textbf{DC component} ($\lambda = 1$): $h(1) = 1.0$ --- preserved exactly.  The total mass is unchanged.
    \item \textbf{Mid-frequency} ($\lambda = 0$): $h(0) = 1-\gamma$ --- moderate attenuation.
    \item \textbf{Highest frequency} ($\lambda = -1$): $h(-1) = 1-2\gamma$ --- strongest attenuation.
\end{itemize}
At $\gamma = 0.26$: the DC component is fully preserved, mid-frequencies are attenuated to 74\%, and the highest-frequency noise is attenuated to 48\%.
This is the spectral signature of a filter that removes oscillatory noise (where adjacent links have wildly different counts due to sampling) while preserving the large-scale spatial structure.

\subsubsection{Variance Preservation}

A key concern with any smoothing is \emph{over-smoothing}: blurring the distribution until all links have similar values, destroying the meaningful variation between high-traffic and low-traffic roads.
We measure the fraction of original variance retained as a function of $\gamma$:
\begin{equation}
    \text{VR}(\gamma) = \frac{\text{Var}(c_{\text{smooth}})}{\text{Var}(\tilde{c})}
\end{equation}

Across 49 random timeframes, the variance retained at $\gamma = 0.26$ is \textbf{67\%} (Figure~\ref{fig:variance}).
This means two-thirds of the original distributional structure is preserved---the Gini coefficient decreases by less than 1\%, and the entropy increases by less than 1\%.
The distribution remains highly concentrated on high-traffic corridors; it is not blurred toward uniformity.

\begin{figure}[H]
\centering
\includegraphics[width=\linewidth]{fig_smoothing_variance.png}
\caption{\textbf{Variance preservation under smoothing.} \textbf{(A)}~Fraction of original variance retained vs $\gamma$. At $\gamma = 0.26$ (dashed), 67\% is retained. \textbf{(B)}~Gini coefficient (distribution concentration) barely changes. \textbf{(C)}~Normalised entropy stays well below 1.0 (uniform), confirming the distribution remains structured. \textbf{(D)}~Spectral filter response: the DC component (green) is preserved at 1.0 while the highest frequency (red) is attenuated by 48\%.}
\label{fig:variance}
\end{figure}

\subsubsection{Spatial Correlation Is Real}

To verify that graph diffusion recovers genuine traffic patterns rather than introducing artifacts, we perform \emph{leave-link-out cross-validation}:
for each observed link, remove it from the count vector, smooth the remainder, and check whether the smoothed value at the removed link approximates the original observation.

Across 15 timeframes and 200 held-out links each, smoothing reduces the recovery error monotonically with $\gamma$: from 98.7\% at $\gamma = 0$ (no information) to 92.9\% at $\gamma = 0.60$.
While the absolute recovery is modest (reflecting the inherent difficulty of predicting one link from its sparse neighbourhood), the \emph{consistent improvement} proves that adjacent links carry real predictive information---the spatial correlation exploited by graph diffusion is genuine, not an artifact.

\subsubsection{Selecting $\gamma$}

The optimal $\gamma$ balances two opposing effects:
\begin{itemize}
    \item \textbf{On observed links} (3\% of network): smoothing adds bias by mixing in neighbour values. On these links, $\gamma = 0$ gives the lowest error.
    \item \textbf{On unobserved links} (97\% of network): smoothing is the \emph{only} source of information. On these links, higher $\gamma$ always helps.
\end{itemize}
Since unobserved links have no ground truth, the combined optimum cannot be determined purely from the data.
We therefore calibrate $\gamma$ via cross-validation on the downstream task, which is standard practice in machine learning (analogous to tuning regularisation strength on a validation set).

Split-half CV (49 timeframes, seed~42) measuring both predictive MSE and downstream PageRank MRE identifies $\gamma = 0.26$ as optimal, within a broad plateau from 0.22 to 0.30 where results are insensitive to the exact value (Figure~\ref{fig:smoothing}).
After smoothing, the matrix becomes $2{,}880 \times 99{,}716$ with 100\% density---every link has a positive probability estimate.

\begin{figure}[H]
\centering
\includegraphics[width=\linewidth]{fig6_smoothing_tradeoff.png}
\caption{\textbf{Smoothing parameter cross-validation.} Split-half CV over 49 random timeframes measuring predictive MSE and downstream PageRank MRE as a function of $\gamma$. Both metrics are minimised at $\gamma = 0.26$ (vertical dashed line), with a broad plateau from 0.22 to 0.30.  The optimal $\gamma$ retains 67\% of original variance while achieving 100\% link coverage.}
\label{fig:smoothing}
\end{figure}

\subsubsection{Summary of Smoothing Justification}

\begin{enumerate}
    \item \textbf{Necessity}: 97\% of links are unobserved in any window.  Without imputation, network-wide analysis is impossible.
    \item \textbf{Method}: Graph diffusion uses only road adjacency---a principled spatial prior.  Leave-link-out CV confirms adjacent links carry real predictive signal.
    \item \textbf{Preservation}: At $\gamma = 0.26$, 67\% of variance is retained.  The Gini coefficient and entropy change by $<1\%$---the distribution remains structured, not blurred.
    \item \textbf{Spectral guarantee}: The DC component is preserved exactly; only high-frequency noise is attenuated (by 48\%).
    \item \textbf{Robustness}: The optimal $\gamma$ sits in a plateau (0.22--0.30) where results are insensitive.  The exact value is calibrated via downstream CV, analogous to regularisation tuning.
\end{enumerate}


% ================================================================
\section{Two-Phase PageRank Model}
\label{sec:model}

\subsection{Standard PageRank}
The classical PageRank~\cite{page1999pagerank, brin1998anatomy} computes the stationary distribution of a random walk on a graph:
\begin{equation}
    v^{(k+1)} = d \cdot P^\top v^{(k)} + (1-d) \cdot E
\end{equation}
where $d$ is the damping factor, $P$ is a row-stochastic transition matrix, and $E$ is a teleportation vector (uniform in the original formulation, or empirical in our case).

\subsection{Travel-Time Self-Loops}
\label{sec:selfloops}

In a road network, vehicles spend time \emph{on} each link proportional to its traversal time $\ell_i / s_i$.
Standard PageRank treats all transitions as instantaneous, assigning equal weight to short connectors and long highways.
We introduce self-loops that model dwell time:
\begin{equation}
    \boxed{sw_i = \mu \cdot \frac{\ell_i}{s_i}}
    \label{eq:selfloop}
\end{equation}
where $\mu > 0$ is the dwell-time scaling parameter.
The transition probability from link~$i$ to successor~$j$ becomes:
\begin{equation}
    P[i, j] = \frac{w_{ij}}{sw_i + \sum_k w_{ik}}, \qquad
    P[i, i] = \frac{sw_i}{sw_i + \sum_k w_{ik}}
    \label{eq:transition}
\end{equation}

For a typical highway link ($\ell = 500$\,m, $s = 25$\,m/s, 3~successors):
\[
sw = 20 \times 500/25 = 400, \quad P[i,i] = \frac{400}{400+3} = 99.3\%
\]
After $k$ iterations, the retained mass is $P[i,i]^k$.
This naturally allocates more PageRank mass to links with longer traversal times---which are precisely the high-volume corridors that carry the most traffic.

Figure~\ref{fig:selfloop} shows how the self-loop probability varies with link traversal time across different $\mu$ values. At $\mu = 20$, links with traversal time above 10~seconds have self-loop probabilities exceeding 98\%, while short connectors ($<1$\,s) retain below 90\%. This creates a spectrum of retention strengths that mirrors the physical hierarchy of the road network.

\begin{figure}[H]
\centering
\includegraphics[width=\linewidth]{fig1_selfloop_prob.png}
\caption{\textbf{Self-loop probability vs.\ traversal time.} Each curve corresponds to a different $\mu$ value (5, 10, 15, 20). The self-loop probability $P[i,i] = sw_i / (sw_i + \sum_k w_{ik})$ increases with traversal time $\ell_i/s_i$. At $\mu = 20$ (red), a highway link with 20\,s traversal time has $P[i,i] = 99.3\%$, meaning only 0.7\% of its mass escapes per iteration.}
\label{fig:selfloop}
\end{figure}

The cumulative effect of self-loops over multiple iterations is shown in Figure~\ref{fig:retention}. After 100 iterations, a highway link retains approximately 50\% of its original mass, while a short connector retains less than 5\%. This differential retention is what produces physically realistic traffic distributions.

\begin{figure}[H]
\centering
\includegraphics[width=\linewidth]{fig2_mass_retention.png}
\caption{\textbf{Mass retention over iterations.} Fraction of initial mass remaining on a link after $k$ power iteration steps, for three representative link types at $\mu = 20$: highway ($\ell/s = 20$\,s, blue), arterial ($\ell/s = 8$\,s, green), and connector ($\ell/s = 1$\,s, red). Highway links retain mass 10$\times$ longer, naturally concentrating PageRank on high-volume corridors.}
\label{fig:retention}
\end{figure}

\subsection{Two-Phase Block Structure}

We decompose the random walk into two phases:
\begin{itemize}
    \item \textbf{Up phase} (active driving): follows the directed adjacency with travel-time self-loops.
    \item \textbf{Down phase} (local dispersal): vehicles that have ``arrived'' disperse locally.
\end{itemize}
The $2N \times 2N$ transition matrix is:
\begin{equation}
    \boxed{\MTN = \begin{bmatrix} (1-\beta)\,P_{\text{up}} & \beta\,I_N \\ \mathbf{0} & P_{\text{down}} \end{bmatrix}}
    \label{eq:block}
\end{equation}
where $\beta \in (0,1)$ is the phase-drop rate.
At each step, a fraction $\beta$ of up-phase mass transitions to the down phase.
The optimal $\beta = 0.9$ implies the average trip has approximately $1/\beta \approx 1.1$ forward hops before dispersal.

\subsection{Power Iteration}

The full update rule is:
\begin{equation}
    \boxed{v^{(k+1)} = d \cdot \MTN^\top \, v^{(k)} + (1-d) \cdot \ETN}
    \label{eq:power}
\end{equation}
where $\ETN = [E_N; \,\mathbf{0}]$ injects the smoothed ground truth into the up phase only, and $d = 0.80$ is the damping factor.
Convergence is measured by $\|v^{(k+1)} - v^{(k)}\|_1 < 10^{-6}$ (typically 37--38 iterations).

The final distribution is obtained by collapsing phases:
\begin{equation}
    \vfinal = \frac{v_{\text{up}} + v_{\text{down}}}{\|v_{\text{up}} + v_{\text{down}}\|_1}
\end{equation}


% ================================================================
\section{Demand-Weighted Transitions}
\label{sec:demand}

\subsection{Motivation}
\label{sec:demand_motivation}

At $d = 0.80$, 20\% of the prediction comes from ground truth.
To reduce this dependence, we increase~$d$---but the topology-only matrix does not encode traffic demand patterns, so accuracy degrades rapidly (Table~\ref{tab:damping}).

A grid search over 576 parameter combinations ($\mu, \alpha_s, \alpha_l, \beta$) found \emph{zero feasible solutions} for $d > 0.90$ with MRE~$\leq 0.05$.

\subsection{The $\gamma$ Parameter}

We introduce a single parameter $\gamma$ that weights each transition by the destination link's average demand:
\begin{equation}
    \boxed{P(i \to j) \propto w_{ij} \cdot \text{demand}[j]^\gamma}
    \label{eq:gamma}
\end{equation}
where $\text{demand}[j] = \frac{1}{T}\sum_{t=1}^T E_N^{(t)}[j]$ is the time-averaged popularity of link~$j$ across a training set of $T$ timeframes.
This teaches the matrix that high-traffic links (e.g., I-15) should attract more flow than residential streets.

\subsection{Validation}

We validate through three independent train/test splits:
\begin{enumerate}
    \item \textbf{Temporal}: Train on Sept~1--15, test on Sept~16--30.
    \item \textbf{Day-of-week}: Train on Mon/Wed/Fri, test on Tue/Thu/Sat/Sun.
    \item \textbf{5-fold random CV}: 80\% train, 20\% test, repeated 5 times.
\end{enumerate}

Results (Table~\ref{tab:demand_validation}) confirm the improvement generalises: at $d = 0.96$, $\gamma = 0.20$, the average Overall~MRE across all splits is $\mathbf{0.0385}$---better than the original model at $d = 0.80$ (0.0466), while using only 4\% ground truth.

\begin{table}[t]
\centering
\caption{Overall MRE across three validation strategies.}
\label{tab:demand_validation}
\small
\begin{tabular}{lcccc}
\toprule
Configuration & Temporal & DoW & 5-Fold & \textbf{Avg} \\
\midrule
$d{=}0.80,\;\gamma{=}0$ & 0.0478 & 0.0457 & 0.0464 & 0.0466 \\
$d{=}0.91,\;\gamma{=}0.10$ & 0.0472 & 0.0452 & 0.0459 & \textbf{0.0461} \\
$d{=}0.94,\;\gamma{=}0.15$ & 0.0439 & 0.0420 & 0.0427 & \textbf{0.0429} \\
$d{=}0.96,\;\gamma{=}0.20$ & 0.0395 & 0.0377 & 0.0383 & \textbf{0.0385} \\
\bottomrule
\end{tabular}
\end{table}


% ================================================================
\section{Results}
\label{sec:results}

\subsection{Evaluation Protocol}

We evaluate on 49 random timeframes (seed~42) held out from parameter tuning.
The ground truth is the smoothed popularity vector $E_N$ for each timeframe.
The prediction is $\vfinal$ from the two-phase PageRank.
The primary metric is Mean Relative Error:
\begin{equation}
    \MRE = \frac{1}{|S|}\sum_{i \in S} \frac{|E_N[i] - \vfinal[i]|}{E_N[i] + \epsilon}
\end{equation}
computed over all links (Overall) or the top-100 busiest links (Top-100).

\subsection{Main Results}

\begin{table}[t]
\centering
\caption{Evaluation metrics (49 timeframes, $\mu{=}20$, $\beta{=}0.9$, $d{=}0.80$).}
\label{tab:main_results}
\small
\begin{tabular}{lrr}
\toprule
Metric & Mean $\pm$ Std & Range \\
\midrule
Top-100 MRE & $0.0312 \pm 0.0122$ & 0.020--0.048 \\
Overall MRE & $0.0464 \pm 0.0056$ & 0.041--0.057 \\
Pearson $r$ & $0.9969 \pm 0.0014$ & 0.995--0.998 \\
Rank overlap /100 & $94.4 \pm 4.5$ & 86--99 \\
\bottomrule
\end{tabular}
\end{table}

Table~\ref{tab:main_results} shows the model achieves Top-100~MRE of 3.1\%, Pearson correlation 0.997, and correctly identifies 94 out of 100 busiest links on average.

The sensitivity of MRE to the dwell-time parameter $\mu$ is shown in Figure~\ref{fig:mre_mu}. Without self-loops ($\mu = 0$), the Top-100 MRE is 0.73---the model severely underpredicts traffic on the busiest roads. Increasing $\mu$ yields monotonic improvement: at $\mu = 20$, Top-100 MRE drops to 0.031, a 95.7\% reduction. This is the single most impactful parameter in the model.

\begin{figure}[H]
\centering
\includegraphics[width=\linewidth]{fig4_mre_vs_mu.png}
\caption{\textbf{MRE as a function of $\mu$.} Top-100 MRE (red, left axis) and Overall MRE (blue, right axis) both decrease monotonically with increasing $\mu$. The improvement from $\mu{=}0$ to $\mu{=}20$ is 95.7\% for Top-100 and 76.2\% for Overall. Diminishing returns appear beyond $\mu = 20$.}
\label{fig:mre_mu}
\end{figure}

Figure~\ref{fig:heatmap} shows the joint sensitivity to $\mu$ and $\beta$. The optimal region ($\mu \geq 15$, $\beta \geq 0.7$) is broad and stable, indicating the model is not sensitive to small parameter perturbations.

\begin{figure}[H]
\centering
\includegraphics[width=\linewidth]{fig3_mre_heatmap.png}
\caption{\textbf{Parameter sensitivity heatmap.} Top-100 MRE over a grid of $\mu$ (dwell-time scale) and $\beta$ (phase-drop rate) values, with $d = 0.80$ fixed. The darkest region (lowest MRE, $<0.04$) spans $\mu \in [15, 20]$ and $\beta \in [0.7, 0.95]$. The optimal configuration ($\mu = 20$, $\beta = 0.9$) lies in the centre of this plateau.}
\label{fig:heatmap}
\end{figure}

\subsection{Effect of Damping Factor}

\begin{table}[t]
\centering
\caption{Damping factor sweep (49 timeframes, $\gamma{=}0$).}
\label{tab:damping}
\small
\begin{tabular}{cccc}
\toprule
$d$ & Top-100 MRE & Overall MRE & Pearson $r$ \\
\midrule
0.80 & 0.0312 & 0.0464 & 0.9969 \\
0.84 & 0.0409 & 0.0605 & 0.9949 \\
0.88 & 0.0558 & 0.0817 & 0.9912 \\
0.92 & 0.0820 & 0.1177 & 0.9829 \\
0.96 & 0.1443 & 0.1947 & 0.9565 \\
\bottomrule
\end{tabular}
\end{table}

Without demand weighting ($\gamma = 0$), MRE increases monotonically with $d$ (Table~\ref{tab:damping}).
With $\gamma = 0.20$, the model achieves Overall~MRE~$= 0.0383$ at $d = 0.96$ (Table~\ref{tab:demand_validation})---a 5$\times$ reduction in ground-truth dependence with improved accuracy.

The full optimisation landscape is shown in Figure~\ref{fig:damping}. Panel~(A) demonstrates the baseline problem: without demand weighting, MRE rises steeply with $d$. Panel~(B) shows how even small $\gamma$ values dramatically reduce MRE at high damping, with the green target zone (MRE~$\leq 0.05$) reachable at $d = 0.96$ for $\gamma \geq 0.10$. Panel~(C) directly compares the key configurations.

\begin{figure}[H]
\centering
\includegraphics[width=\linewidth]{fig_damping_final.png}
\caption{\textbf{Damping factor optimisation.} \textbf{(A)}~Baseline sweep ($\gamma{=}0$): Top-100 (red) and Overall (blue) MRE with error bars across 49 timeframes. MRE rises monotonically; the 0.05 target (dashed) is only met at $d{=}0.80$. \textbf{(B)}~With demand weighting: each line is a damping value; increasing $\gamma$ drives MRE below the target even at $d{=}0.96$. Green shading marks the feasible zone. \textbf{(C)}~Bar chart comparing the original ($d{=}0.80$, red highlight) against three optimised configurations.}
\label{fig:damping}
\end{figure}

To verify that this improvement is not an artifact of data leakage, we validate across three independent train/test splits (Figure~\ref{fig:validation}). In each case, the demand prior is computed only from the training set and the model is evaluated on held-out data it has never seen. The results are consistent across all three strategies, confirming the demand prior captures a stable structural property of the network.

\begin{figure}[H]
\centering
\includegraphics[width=\linewidth]{fig_validation.png}
\caption{\textbf{Rigorous validation across three data splits.} Each panel shows 8 configurations evaluated on a different held-out test set. \textbf{Left:}~Temporal split (train Sept~1--15, test Sept~16--30). \textbf{Centre:}~Day-of-week split (train MWF, test TThSSu). \textbf{Right:}~5-fold random CV (80/20 split, averaged over 5 folds). Solid bars: Overall MRE; transparent bars: Top-100 MRE. The demand-weighted configurations ($\gamma{=}0.15$--$0.20$) consistently meet the 0.05 target at $d{=}0.94$--$0.96$.}
\label{fig:validation}
\end{figure}


% ================================================================
\section{Event Detection and Prediction}
\label{sec:events}

We validate the model against two documented real-world events:
an unplanned traffic crash and a planned sporting event.

\subsection{I-15 Multi-Vehicle Crash (September~20, 2018)}
\label{sec:crash}

A four-vehicle collision blocked multiple southbound lanes of I-15 at 14400~South, Draper, from approximately 7:30--11:00~AM~\cite{gephardtdaily2018}.
We compare crash-day PageRank mass on the 21~I-15~SB links against three baseline Thursdays (Sept~6, 13, 27).

\begin{table}[t]
\centering
\caption{I-15 crash: PageRank deviation from baseline.}
\label{tab:crash}
\small
\begin{tabular}{llrrr}
\toprule
Time & Phase & I-15 SB & I-15 NB & Surface \\
\midrule
06:00 & Pre-crash & $+1.7\%$ & $+96.5\%$ & $+10.4\%$ \\
08:30 & Peak & $\mathbf{-18.9\%}$ & $\mathbf{-92.7\%}$ & $-58.6\%$ \\
10:30 & Rebound & $\mathbf{+291.4\%}$ & $-40.9\%$ & $-8.3\%$ \\
11:30 & Recovery & $+36.7\%$ & $+37.1\%$ & $+15.2\%$ \\
\bottomrule
\end{tabular}
\end{table}

The model captures the complete crash lifecycle (Table~\ref{tab:crash}):
pre-crash calibration ($+1.7\%$, confirming fair comparison), peak disruption ($-18.9\%$~SB, $-92.7\%$~NB from rubbernecking), clearance rebound ($+291.4\%$~SB as pent-up traffic releases), and recovery ($+36.7\%$).
The demand-weighted model ($d = 0.96$, $\gamma = 0.20$) produces virtually identical deviations ($-18.9\%$, $+292.4\%$), confirming crash detection with only 4\% ground truth.

The complete crash analysis is summarised in Figure~\ref{fig:crash_panel}. Panel~(A) shows the smoothed I-15~SB traffic throughout the day, with the crash day (red) deviating from the baseline band during 08:00--11:00. Panel~(B) shows the raw GPS counts---extremely sparse (0--5 per 15-minute window), illustrating why smoothing and PageRank are necessary. Panel~(C) plots the percentage deviation from baseline for all three corridor categories, revealing the distinctive ``dip-then-spike'' crash signature. Panel~(D) shows surface streets, where suppressed traffic suggests area-wide congestion rather than simple freeway-to-arterial diversion.

\begin{figure}[H]
\centering
\includegraphics[width=\linewidth]{fig5_combined_panel.png}
\caption{\textbf{I-15 crash analysis: four-panel summary.} \textbf{(A)}~Smoothed I-15~SB traffic mass. Red: crash day (Sept~20); blue: baseline average $\pm 1\sigma$ from three Thursdays. The crash window (7:30--11:00~AM, shaded red) shows clear deviation. \textbf{(B)}~Raw GPS counts on I-15~SB: 0--5 per 15-min slot, illustrating the sparsity that motivates our pipeline. \textbf{(C)}~PageRank deviation from baseline for SB (red, $-18.9\%$ at 08:30, $+291\%$ at 10:30), NB (orange, $-92.7\%$ rubbernecking), and surface (green). \textbf{(D)}~Surface arterials show suppressed traffic during the crash window.}
\label{fig:crash_panel}
\end{figure}

\subsection{Utah vs.\ \#10 Washington Football Game (September~15, 2018)}
\label{sec:game}

A near-capacity crowd of 47,445 attended the Saturday night game at Rice-Eccles Stadium (kickoff 8:00~PM)~\cite{utahutes2018}.
Documented road closures include 500~South lane reduction and Guardsman~Way full closure~\cite{utahathletics2018}.
We compare game day against four baseline Saturdays (Sept~1, 8, 22, 29) on 93~event-affected links.

The model detects the full game-day cycle:
$-85.9\%$ on 500~South at 14:00 (road closures),
$+463.4\%$ at 19:00 (arrival surge),
$+247.9\%$ at stadium core at 23:00 (post-game departure).

Figure~\ref{fig:game} summarises the game day analysis across four panels. The smoothed timeseries (A) shows traffic on the 93 event-affected links surging well above the baseline band during the 18:00--20:00 arrival window. The deviation plot (B) reveals the complete event signature: negative deviations during daytime road closures, a massive positive spike at arrival, negative during the game itself (everyone is inside), and a second positive spike at departure. The $\mu$-control prediction (C) shows that reducing $\mu$ on event links brings prediction error from 200--300\% down to 2--15\% during closure phases. Panel~(D) traces the optimal $\mu$ over time: near-zero during closures, returning to 20 during arrival surges where mass \emph{increases}.

\begin{figure}[H]
\centering
\includegraphics[width=\linewidth]{fig_game_summary.png}
\caption{\textbf{Utah vs.\ \#10 Washington game day analysis.} \textbf{(A)}~Smoothed traffic on 93 event-affected links near Rice-Eccles Stadium. Red: game day; blue: baseline average from 4~Saturdays. The 19:00 arrival surge peaks at 3$\times$ baseline. \textbf{(B)}~PageRank deviation for Event Links (blue), 500~South (orange), and Guardsman Way (purple). Shaded regions mark pre-game, game, and post-game phases. \textbf{(C)}~$\mu$-control prediction: red bars show error at $\mu{=}20$ (no conditioning); green bars show error at optimal $\mu$. Improvement exceeds 90\% at 14:00--18:00. \textbf{(D)}~Optimal $\mu$ by time slot: near-zero during road closures (14:00--18:00, 21:00), returning to 20 during arrival (19:00) and departure (23:00) surges.}
\label{fig:game}
\end{figure}


% ================================================================
\section{The Role of $\mu$: Per-Link Disruption Control}
\label{sec:mu}

\subsection{Why $\mu$ Is the Only Effective Control Parameter}

The self-loop probability $P[i,i] = sw_i / (sw_i + \sum_k w_{ik})$ is dominated by $sw_i = \mu \cdot \ell_i / s_i$ at $\mu = 20$ (typically $> 99\%$).
We tested five types of intervention on crash-affected links:

\begin{table}[t]
\centering
\caption{Effect of different interventions on I-15~SB corridor mass.}
\label{tab:interventions}
\small
\begin{tabular}{lrc}
\toprule
Intervention & What changes & Max effect \\
\midrule
Speed $= 0.1$\,m/s & Outgoing weights & $-0.44\%$ \\
Lanes $= 1$ & Outgoing weights & $-0.54\%$ \\
Remove from adjacency & Zero outgoing & $+0.67\%$ \\
Zero demand weight & Zero attraction & $+0.67\%$ \\
\textbf{$\mu = 0.01$} & \textbf{Self-loop: $99.5\% \to 9.6\%$} & $\mathbf{-64.4\%}$ \\
\bottomrule
\end{tabular}
\end{table}

Only $\mu$ produces meaningful changes (Table~\ref{tab:interventions}) because it operates on the 99.5\% retained mass, while all other parameters affect only the 0.5\% outgoing flow.

\subsection{Crash Prediction Without Crash-Day Data}

Using September~13 (normal Thursday) as input and reducing $\mu$ on the 21~I-15~SB links:

\begin{table}[t]
\centering
\caption{I-15~SB prediction error at key crash times.}
\label{tab:mu_crash}
\small
\begin{tabular}{lccrr}
\toprule
Time & $\mu{=}20$ & Best $\mu$ & Best err & Improv. \\
\midrule
08:00 & 34.1\% & 0.10 & \textbf{0.0\%} & $+100\%$ \\
08:30 & 179.1\% & 0.01 & \textbf{11.5\%} & $+93.6\%$ \\
08:45 & 380.7\% & 0.00 & \textbf{26.4\%} & $+93.1\%$ \\
10:15 & 142.2\% & 0.01 & \textbf{2.6\%} & $+98.1\%$ \\
\bottomrule
\end{tabular}
\end{table}

During the active diversion phase, $\mu$-control reduces corridor error by 91--100\% (Table~\ref{tab:mu_crash}).
The same approach on the football game yields 94--98\% error reduction on stadium-area links during closure phases.

\subsection{Cross-Validated Baseline: Tuesday $\mu$ Profile}

To confirm that $\mu = 20$ is genuinely optimal on normal days, we perform leave-one-out cross-validation across four Tuesdays (no known events).
For each of 96~timeframes, we sweep $\mu$ and find the value minimising Overall~MRE.

\textbf{Result}: The optimal $\mu$ stays in the range 10--20 across all 96~slots, with only 1\% improvement over fixed $\mu = 20$.
The correlation between raw ground-truth gap (between Tuesdays) and model prediction error is $r = 0.9999$---the model adds less than 4\% distortion to the input data.
The model error is 96.6\% of the raw input gap, confirming that the error comes from week-to-week GPS variation, not model miscalibration.

This establishes the baseline: $\mu \approx 20$ on normal days, so any significant deviation signals a real disruption.

The cross-validated $\mu$ profile is shown in Figure~\ref{fig:tuesday}. Panel~(A) shows the median optimal $\mu$ across 4 folds with $\pm 1\sigma$ shading---it stays firmly between 10 and 20 throughout the day. Panel~(B) shows individual fold profiles, demonstrating consistency across all four Tuesdays. Panel~(C) overlays the MRE at $\mu = 20$ (red) against the MRE at optimal $\mu$ (green): the lines nearly overlap, confirming that $\mu = 20$ is already near-optimal on normal days.

\begin{figure}[H]
\centering
\includegraphics[width=\linewidth]{fig_mu_profile_tuesday.png}
\caption{\textbf{Cross-validated Tuesday $\mu$ profile.} Leave-one-out CV across 4 Tuesdays (Sep~4, 11, 18, 25). \textbf{(A)}~Median optimal $\mu$ (blue line) with $\pm 1\sigma$ band across folds. The default $\mu{=}20$ (grey dashed) falls within the confidence band at all 96 time slots. \textbf{(B)}~Individual fold profiles (coloured lines) agree closely. \textbf{(C)}~MRE comparison: fixed $\mu{=}20$ (red) vs per-slot optimal $\mu$ (green). The difference is negligible ($+1.0\%$ average improvement), confirming the model is well-calibrated on normal days.}
\label{fig:tuesday}
\end{figure}

Figure~\ref{fig:gt_model} provides the deepest insight: the relationship between the raw ground-truth gap and the model prediction error. Panel~(A) overlays both quantities across the 24-hour cycle---they follow identical patterns ($r = 0.9999$). Panel~(B) shows the model's ``self-distortion'': the MRE between the input $E_N$ and the PageRank output $\text{PR}(E_N)$ is only 0.039 (3.9\%), confirming that PageRank is a nearly transparent transformation at $d = 0.96$. Panel~(C) brings all three together, showing that the raw input gap (0.264) dominates the model output gap (0.255), with self-distortion (0.039) as a small residual.

\begin{figure}[H]
\centering
\includegraphics[width=\linewidth]{fig_mu_groundtruth_vs_model.png}
\caption{\textbf{Ground truth gap vs model error.} \textbf{(A)}~Red: raw MRE between held-out Tuesday $E_N$ and training average $E_N$ (``input gap''). Blue dashed: model prediction error at $\mu{=}20$ (``output gap''). The curves overlap with $r = 0.9999$---model error is determined by input variation, not model parameters. Per-fold gaps shown as faint coloured lines. \textbf{(B)}~Model self-distortion: MRE($E_N$, PR($E_N$)) is only 0.039. PageRank barely transforms its input at $d{=}0.96$. \textbf{(C)}~All three layers: the input gap (red shaded) is the floor; the output gap (blue) tracks it at 96.6\%; self-distortion (green) is minimal.}
\label{fig:gt_model}
\end{figure}

\subsection{Practical Protocol}

For real-time event detection and scenario analysis:
\begin{enumerate}
    \item Build the demand-weighted matrix ($\gamma = 0.20$) once using historical data.
    \item For \textbf{detection}: run PageRank with current GPS data and compare to baseline. Deviations beyond $\pm 2\sigma$ flag anomalies.
    \item For \textbf{prediction}: when a disruption is reported on specific links, set $\mu \to 0.01$ on those links and rerun PageRank with the most recent normal-day data. The output predicts the new traffic equilibrium in $<$10\,ms on GPU.
    \item \textbf{Severity mapping}: $\mu = 20$ (normal), $\mu \sim 1$ (partial restriction), $\mu \sim 0.01$ (full closure).
\end{enumerate}


% ================================================================
\section{Conclusion}
\label{sec:conclusion}

We presented a Two-Phase PageRank model with travel-time self-loops and demand-weighted transitions for urban traffic distribution estimation.
The model achieves 3.1\% Top-100~MRE on 49 held-out timeframes and detects real-world disruptions (crashes, sporting events) through PageRank deviation analysis.

Three parameters control the model:
\begin{itemize}
    \item \textbf{Damping $d$}: balances matrix topology against ground truth (optimal: 0.96 with demand weighting).
    \item \textbf{Demand sensitivity $\gamma$}: encodes structural traffic demand into the matrix (optimal: 0.20, validated across 3~independent splits).
    \item \textbf{Dwell-time $\mu$}: controls per-link mass retention---the only effective parameter for local disruption modelling (default: 20, disrupted: 0.01).
\end{itemize}

The key insight is the separation of concerns: $\gamma$ handles global prediction accuracy, while $\mu$ handles local disruption response.
Together, they enable a model that uses only 4\% ground truth yet achieves high accuracy and can predict event-driven traffic redistribution without any event-day data.

Future work includes time-varying $\mu$ schedules matched to incident timelines, extension to multi-modal networks, and integration with real-time incident detection systems.

\bibliographystyle{plainnat}
\bibliography{references}

\end{document}
